\section*{الفصل \ref{ch:b2:dipole-radiation} --- إشعاع ثنائي القطب والتبعثر}

\begin{solution}{exo:b2:dipole-radiation:1}
$p_0 = I_0\ell/\omega = \qty{3.2e-9}{C.m}$؛ و $E = \mu_0\omega^2p_0\sin\theta/4\pi r$: \qty{0.13}{V/m}
عند $\theta = 90^\circ$، و \qty{0.063}{V/m} عند \qty{30}{\degree}؛ و $\mathcal P = \mu_0\omega^4p_0^2/12\pi c =
\qty{175}{W}$؛ و $R_{\text{إشعاع}} = 2\mathcal P/I_0^2 = \qty{88}{\ohm}$.
\end{solution}

\begin{solution}{exo:b2:dipole-radiation:2}
$790(\ell/\lambda)^2$: \qty{8.8e-5}{\ohm}، \qty{8.8e-3}{\ohm}، \qty{0.88}{\ohm}، \qty{88}{\ohm}. وأما النحاس:
$0.026$، $0.084$، $0.27$، \qty{0.84}{\ohm} (\omterm{thm:b2:waves-in-media:skin}{بالأثر القشري}): فالمردودات $0.3\%$، $10\%$،
$77\%$، $99\%$. وعند أطوال الموجة الطويلة لا تكون \omterm{ex:b2:dipole-radiation:antenna}{مقاومة الإشعاع} فوق ضياعاته إلا
لصارٍ يبلغ كسرًا معتبرًا من $\lambda$.
\end{solution}

\begin{solution}{exo:b2:dipole-radiation:3}
$400 : 550 : 700 = 9.4 : 2.6 : 1$. و $(10/3)^4 = 120$: فرادار \qty{3}{cm}
يرى القطرات الصغيرة أحسن بمئة مرة؛ ورادار \qty{10}{cm} ينظر
عبر المطر إلى ما وراءه.
\end{solution}

\begin{solution}{exo:b2:dipole-radiation:4}
$\sin^2\chi/(1 + \cos^2\chi)$: $0.14$، $0.60$، $1$، $0.60$. وجّهه إلى \qty{90}{\degree} من
الشمس (أي دائرة كبرى عبر السماء)؛ وحين تكون الشمس في السمت يقع الأفق
كله على \qty{90}{\degree} ويكون الاستقطاب على امتداده. والتبعثر المتعدد،
وانعكاس الأرض، وتباين مناحي الجزيئات تضيف ضوءًا غير
مستقطب: فنحو $75\%$ في أحسن الأحوال.
\end{solution}

\begin{solution}{exo:b2:dipole-radiation:5}
(أ) $\int\sin^2\theta\,\dd\Omega = 8\pi/3$. (ب) $\int_{60^\circ}^{120^\circ}\sin^3\theta\,\dd\theta/
(4/3) = 0.917/1.333 = 69\%$. (ج) $\sin^2\theta = \tfrac12$: $\theta = \qty{45}{\degree}$، أي
حزمة عرضها \qty{90}{\degree}. (د) والعظم $1$ مقابل المتوسط $2/3$: أي $1.5$.
\end{solution}

\begin{solution}{exo:b2:dipole-radiation:6}
(أ) $p/4\pi\varepsilon_0r^3 = \omega^2p/4\pi\varepsilon_0c^2r$ عند $r_0 = c/\omega = \lambda/2\pi$. (ب) \qty{48}{m}:
فعلى \qty{10}{m} الحقل شبه الساكن، وعلى \qty{10}{km} الحقل المشعّ. (ج)
$r_0 = \qty{4.8}{cm}$: فالرأس في المنطقة القريبة. (د) وحقلا المنطقة القريبة $\vect E$
و $\vect B$ في التربيع: فمتوسط شعاع بوينتنغ معدوم،
وتُخزَّن الطاقة وتُعاد، كما في مكثفة.
\end{solution}

\begin{solution}{exo:b2:dipole-radiation:7}
(أ) $8\pi r_e^2/3 = \qty{6.65e-29}{m^2}$. (ب) $n_e\sigma_TL = 10^{14} \times 6.65 \times 10^{-29}
\times 10^9 = 7 \times 10^{-6}$: أي جزء من مليون من الكرة الضوئية، بلا لون لأن
\omterm{prop:b2:dipole-radiation:rayleigh}{تبعثر طومسون} كذلك، ويغرقه ضوء الشمس المتبعثر في السماء
إلا عند الكسوف. (ج) و \qty{10}{keV} $\gg$ \qty{300}{eV}: فتستجيب الإلكترونات
استجابة حرة، $6\sigma_T = \qty{4e-28}{m^2}$. (د) والامتصاص الكهرضوئي ينمو وفق $Z^4$
ويفرز الكلسيوم؛ أما التبعثر فهو نفسه لكل إلكترون في العظم
واللحم.
\end{solution}

\begin{solution}{exo:b2:dipole-radiation:8}
(أ) $\tau = 0.36$، $0.10$، $0.038$: $T = 0.70$، $0.90$، $0.96$. (ب) و $\tau \times 38 = 13.6$،
$3.8$، $1.45$: $T = 1 \times 10^{-6}$، $0.02$، $0.23$؛ والأحمر على الأزرق $\sim 2 \times 10^5$. (ج)
وجوّ الأرض يحني ضوء الغروب إلى داخل الظل، محمرًّا
بالتبعثر نفسه. (د) وقرب الأفق يطول المسار: فيتبعثر
الأزرق خارجًا من الضوء المتبعثر نفسه، ويضيف التبعثر
المتعدد بياضًا.
\end{solution}

\begin{solution}{exo:b2:dipole-radiation:9}
(أ) على استقامة الزوجين يعطي فرق المسير $\lambda/2$
تلاشيًا؛ وعرضًا يعطي تجمّعًا. (ب) وعلى تعاكس الطور: العكس ---
إشعاع على الاستقامة (الرصّة الطرفية)، ولا شيء عرضًا. (ج) ونحو $+x$:
تأخّر المسير $-\pi/2$ زائدًا طور التغذية $+\pi/2$، أي توافق في الطور؛ ونحو $-x$:
$-\pi$، أي تلاشٍ: فيكون النقش قلبيًّا، من جهة واحدة. (د) والصواري تشكّل التغطية
وتحمي الجيران؛ والرصّة الطورية توجّه بتغيير أطوار
التغذية، في ميكروثوانٍ.
\end{solution}

\begin{solution}{exo:b2:dipole-radiation:10}
(أ) $\langle\mathcal P\rangle = e^2\omega_0^4x_0^2/12\pi\varepsilon_0c^3$، و $W = \tfrac12m\omega_0^2x_0^2$:
$\Gamma = \mathcal P/W = e^2\omega_0^2/6\pi\varepsilon_0mc^3$. (ب) و $\omega_0 = \qty{3.2e15}{rad/s}$: $\Gamma =
\qty{6.4e7}{s^{-1}}$ --- وهي القيمة المستعملة. (ج) و \qty{16}{ns}؛ و $\Delta\lambda = \lambda^2\Gamma/2\pi c =
\qty{1.2e-14}{m}$، أي جزء من مئة من البيكومتر (\qty{10}{MHz}). (د) و $Q = 5 \times 10^7$.
\end{solution}

\begin{solution}{exo:b2:dipole-radiation:11}
(أ) \qty{75}{m}؛ و $I_0 = \sqrt{2\mathcal P/R} = \qty{53}{A}$. (ب) و $\langle\Pi\rangle = 1.5 \times
5 \times 10^4/2\pi \times 10^8 = \qty{1.2e-4}{W/m^2}$، و $E_0 = \sqrt{2\Pi/\varepsilon_0c} = \qty{0.30}{V/m}$.
(ج) \qty{0.30}{V}. (د) وتيارات العودة تجري في الأرض؛ فالتربة المقاومة
تضيف ضياعات وتفسد المرآة؛ والأسلاك القطرية المدفونة تعطيها
مسارًا نحاسيًّا.
\end{solution}

\begin{solution}{exo:b2:dipole-radiation:12}
(أ) $p_0 = \alpha E_0$ يعطي $\sigma = \alpha^2\omega^4/6\pi\varepsilon_0^2c^4$؛ ومع $\alpha = \varepsilon_0(n^2
- 1)/N$ و $\omega = 2\pi c/\lambda$: $\sigma = 8\pi^3(n^2 - 1)^2/3N^2\lambda^4$. (ب) و $n^2 - 1 =
5.8 \times 10^{-4}$: $\sigma = \qty{4.9e-31}{m^2}$، والمسار الحر الوسطي $1/N\sigma = \qty{80}{km}$. (ج)
و $8/80 = 0.1$. (د) ويُقاس $\sigma$ من \omterm{def:b2:dispersion-wave-packets:complex}{وهن} السماء، و $n$
معلومة: فينتج $N$، ومن ثَمّ عدد أفوغادرو.
\end{solution}

\begin{solution}{pb:b2:dipole-radiation:1}
\textbf{1.} $\lambda = \qty{300}{m}$، و $h = \qty{75}{m}$؛ و $I_0 = \sqrt{2 \times 10^5/36} = \qty{75}{A}$؛
و $V = RI_0 = \qty{2.7}{kV}$.

\textbf{2.} $\ell_{\text{فعّال}} = \qty{48}{m}$: $p_0 = 75 \times 48/6.3 \times 10^6 = \qty{5.7e-4}{C.m}$.

\textbf{3.} $\langle\Pi\rangle = 1.5 \times 10^5/2\pi r^2$: \qty{2.4e-4}{W/m^2} و \qty{2.4e-6}{W/m^2}؛
و $E_0 = \qty{0.42}{V/m}$ و \qty{0.042}{V/m}.

\textbf{4.} السوط: \qty{42}{mV}. والقضيب: $N\mu_{\text{فعّال}}A\omega B_0 = 10^4 \times 10^{-4} \times 6.3 \times
10^6 \times 1.4 \times 10^{-10} = \qty{0.9}{mV}$ --- أي أقلّ، لكنه صغير وانتقائي.

\textbf{5.} $A = 1.5\lambda^2/4\pi = \qty{1.1e4}{m^2}$: $P = \Pi A = \qty{26}{mW}$.

\textbf{6.} $36/38 = 95\%$؛ و \qty{5}{kW} حرارةً.

\textbf{7.} $r/\lambda = 33$ على \qty{10}{km}: أي \omterm{thm:b2:dipole-radiation:field}{منطقة الإشعاع}؛ وعلى \qty{100}{m}،
$r \approx \lambda/3$: أي منطقة الانتقال، حيث لا يزال الحقل شبه الساكن
يُحسب له حساب.

\textbf{8.} نهارًا تمتص الطبقة D الموجة التي كانت لولاها
تبلغ الطبقات العاكسة؛ وليلًا تزول فتحمل الموجة السماوية
المحطة مئات الكيلومترات.

\textbf{9.} $\sigma(450) = \qty{1.0e-30}{m^2}$، و $\sigma(650) = \qty{2.4e-31}{m^2}$: والنسبة $4.4$.

\textbf{10.} المسارات الحرة الوسطى $39$، $87$، \qty{170}{km}؛ و $\tau = 0.21$، $0.09$، $0.05$:
أي $19\%$، $9\%$، $5\%$ متبعثرة.

\textbf{11.} $N\sigma I = 2.5 \times 10^{25} \times 4.6 \times 10^{-31} \times 300 = \qty{3.5}{mW/m^3}$،
في $4\pi$ بالنقش $(1 + \cos^2\chi)$ من أجل \omterm{def:b2:plane-waves-polarization:states}{الضوء الطبيعي}.

\textbf{12.} يبعثر العمود $3.5 \times 10^{-3} \times 8000 = \qty{28}{W}$ من
\qty{300}{W} ($9\%$)؛ وبتوزّعها على نصف الكرة تعطي السماء نحو
\qty{4}{W.m^{-2}.sr^{-1}}، مقابل $300/6.8 \times 10^{-5} = \qty{4e6}{W.m^{-2}.sr^{-1}}$ للشمس:
أي جزء من مليون في الستراديان --- وهي الرتبة الصحيحة.

\textbf{13.} مستقطب تمامًا عند \qty{90}{\degree} في التبعثر الأحادي؛ والتبعثر
المتعدد والأرض يخفضانه إلى نحو $75\%$؛ والنحل والنمل يقرأان
النقش ليجدا الشمس خلف الغيوم.

\textbf{14.} $\tau \times 38$: الأزرق $8$ ($T = 3 \times 10^{-4}$)، والأحمر $1.8$ ($T = 0.17$): أي شمس
حمراء غامقة؛ وفوق الرأس يصل الضوء بعد مسار قصير فيكون أزرق.

\textbf{15.} لا جوّ، فلا شيء يبعثر: فالسماء سوداء بجوار
الشمس.

\textbf{16.} غبار الميكرومتر يبعثر تبعثر جسيمات مي: أمامًا
وأحمر، فيعطي سماء بلون الكرمل؛ وقرب الشمس الغاربة يكون
الضوء المتبعثر أمامًا أزرق.

\textbf{17.} $8\pi r_e^2/3 = \qty{6.65e-29}{m^2}$.

\textbf{18.} $a = v^2/2d = \qty{6.4e21}{m/s^2}$؛ و $\mathcal P = e^2a^2/6\pi\varepsilon_0c^3 =
\qty{2e-10}{W}$ خلال $2d/v = \qty{2.5e-14}{s}$: أي \qty{6e-24}{J}، و $10^{-10}$ من
\qty{100}{keV} --- والمردود الحقيقي قرب $1\%$، لأن الكبح يقع
في لقاءات قريبة عنيفة مع النوى، لا في تباطؤ منتظم
لطيف.

\textbf{19.} $hc/E = \qty{12.4}{pm}$.

\textbf{20.} \qty{2}{kW} في الحزمة، ونحو \qty{20}{W} أشعةً سينية؛ وكيلوواطان
حرارةً في بقعة --- ولذلك يدور المصعد لينشرها.

\textbf{21.} $a = v^2/r = \qty{9e22}{m/s^2}$، و $\mathcal P = \qty{5e-8}{W}$؛ وتذهب \qty{13.6}{eV}
في \qty{5e-11}{s}.

\textbf{22.} الشحن الكلاسيكي في حركة مرتبطة يشعّ باستمرار؛
أما الميكانيك الكمومي فله حالات مستقرة لا تشعّ، وأدناها
لا مهوى تحتها.

\textbf{23.} $a = c^2/R = \qty{9e14}{m/s^2}$: أي \qty{5e-24}{W} بلارمور، و $\gamma^4 \sim 10^{16}$
ضعف ذلك في الواقع: فعشرات الكيلوواطات من الضوء السنكروتروني من
حزمة مخزونة.

\textbf{24.} $\sigma_TI = \qty{6.7e-26}{W}$؛ و $1.5 \times 10^{25}$ إلكترون من أجل واط.

\textbf{25.} الصاري: \qty{5.7e-4}{C.m} عند \qty{1}{MHz}، و \qty{100}{kW}، وكعكة
شاقولية. والجزيء: $\sim$\qty{1e-35}{C.m} عند \qty{500}{THz}، و \qty{1e-28}{W}،
وكعكة حول الحقل. والإلكترون الحر: طومسون، \qty{1e-25}{W} في
ضوء الشمس، والنقش نفسه.
\end{solution}
